% Thesis abstract in English (follows the Vietnamese abstract)
\begin{otherlanguage}{english}
\chapter*{ABSTRACT}
\addcontentsline{toc}{chapter}{Abstract}

This thesis studies the problem of forecasting stock price movement trends on the
Vietnamese stock market by combining Vietnamese financial news with price
time-series data. The motivation stems from the fact that stock prices reflect
not only trading history but also react to information about business results,
investment plans, changes in corporate governance, policies, legal risks, and
unusual events. The scope is limited to the HOSE exchange and a group of VN30
constituents. The objective is to determine the direction of the next trading
session across three classes -- up, sideways, and down -- rather than to forecast
the absolute price value directly.

The system is organized into two branches. The first branch uses PhoBERT, a
pre-trained Transformer model for Vietnamese, to classify financial news into
three classes: positive, neutral, and negative \cite{Nguyen2020_200300744v3}.
Labels are interpreted as the expected impact of the content on the related firm
or stock rather than the general tone of the text. After inference, the system
retains the probabilities of all three classes instead of only the
highest-probability label. These probabilities are aggregated by ticker and
trading day together with the article count, class proportions, dispersion, and
a flag indicating whether news is present.

The second branch uses adjusted OHLCV data together with causal features
comprising return, log return, the ratio of closing price to moving average,
volume change, and volatility. The two branches are aligned by ticker and date
after applying point-in-time rules. News is converted to the Vietnam time zone;
articles published after the information cut-off, on weekends, or on holidays are
mapped to the next trading session. This rule ensures that the timestamp of each
article is consistent with the session that uses it; the point-in-time limitation
of the retrospective sentiment checkpoint is stated explicitly in the results
chapter.

Regarding architecture, the models are arranged in tiers to separate the question
about data from the question about complexity. The first tier consists of a model
that always predicts the most frequent class and a model that samples from the
class distribution. The baseline tier comprises a price-only LSTM and a
two-branch LSTM combining price with sentiment. The temporal fusion transformer
is only an extension option when data, resources, and baseline results permit.

Methodologically, the thesis builds a pipeline covering literature survey, data
collection and cleaning, duplicate removal, provenance logging, news-to-ticker
mapping, session alignment, sentiment model fine-tuning, feature engineering,
training, and evaluation. The news dataset includes the title, the content when
collectable, the source, the URL, the publication timestamp, and the related
tickers. Articles mentioning multiple tickers are stored under a many-to-many
relation; records missing a timestamp or without an identifiable ticker are
flagged and are not entered directly into the main feature table.

Price data are split chronologically by walk-forward validation with five
disjoint test folds; there is no separate final test set because the
configuration is fixed before the run and each test fold is scored only once. The
scaler, imputer, label thresholds, feature selection, and price-branch parameters
are learned only from the training portion of each window. The
sentiment-augmented model is compared with the price-only model on the same
training, validation, and test days, the same preprocessing, the same random seed
or seed set, and the same stopping criterion. The control condition keeps the
architecture and the news coverage unchanged in order to isolate the
retrospective gap due to sentiment probabilities from the effect of architecture
or sample splitting.

Regarding data, the news repository holds 10,695 records with URL, publication
timestamp, and content within the 2020--2025 window, combined with 12,624
news--ticker links across 10 VN30 tickers and 14,990 OHLCV price rows for the
2020--2025 period. The curated in-domain sentiment label set comprises 1,306
samples, alongside a CafeF seed set of 999 samples. Stratified five-fold
cross-validation on the 1,306 samples attains macro F1 0.7298, accuracy 0.8125,
and balanced accuracy 0.7689; the five out-of-fold test portions together contain
59 NEGATIVE samples with a coverage of 0.6803.

After the configuration was frozen, a deployment checkpoint was fine-tuned on all
1,306 labels for five fixed epochs and re-scored on 14,256 news--ticker links, of
which 12,624 fall inside the 2020--2025 window and 12,613 anchor to a session
inside the window. The forecasting table comprises 14,990 ticker--session rows,
of which 6,948 rows have news. On the forecasting task, the price-only LSTM
attains macro F1 0.3756, balanced accuracy 0.3915, and one-vs-rest AUC-ROC
0.5714; this macro F1 is above the random level 0.3289 and the majority-class
level 0.1453. The minimal TFT tier attains 0.3619, below the LSTM. Because the
three-class labels are cut at tertiles, the classes are nearly balanced by design
and the random level of the task is 0.333.

The standalone contribution of the sentiment signal is measured by a control that
keeps the two-branch architecture, the same articles, timestamps, and article
counts, but replaces the true sentiment probabilities with a neutral prior. In
the retrospective run with the checkpoint fine-tuned on all 1,306 labels, the
direct difference is $+0.0113$ macro F1, of which the effect of architecture,
news presence, and news volume is $+0.0034$ and the window-level information
contribution is $+0.0080$; the 95\% day-level confidence interval $[-0.0025;
+0.0186]$ contains 0.

Because the retrospective checkpoint learns from labels published after the test
days, the figures above are retrospective differences only. The out-of-sample run
--- a frozen checkpoint trained only on the 1,049 labels published before the
first test observation date, together with a size- and class-matched control
checkpoint to isolate the effect of training-set size --- has been measured and
verified. The out-of-sample information contribution of the two-branch LSTM is
$+0.0094$ macro F1 at window level, with a day-level confidence interval
$[-0.0009; +0.0196]$ containing 0; the size-matched control (seed 7) yields
$+0.0028$ at window level, with a day-level confidence interval $[-0.0074;
+0.0133]$ containing 0. All three point estimates are positive and lie between
$+0.003$ and $+0.009$, yet every confidence interval contains 0; moreover, these
runs constitute multiple comparisons on the same dataset, so there is no basis to
claim statistical significance. Four supplementary analyses were carried out to
probe the robustness of this finding: rank correlation between sentiment scores
and returns without a model, a label set built on excess return relative to
same-session peers, a forecast-horizon ablation, and an economic diagnostic under
an exploratory, non-tradable framing. None of the four finds a reproducible
improvement. The conclusion within the measured scope remains that there is not
yet sufficiently strong evidence that sentiment improves forecasting; the data
also do not rule out a small improvement. The results must not be interpreted as
a causal relationship or as investment advice.

\vspace{0.5cm}
\noindent Keywords: stock price trend forecasting, sentiment analysis, PhoBERT,
Transformer, time series, HOSE, VN30.
\end{otherlanguage}
